\ifdefined\gkver\else\def\gkver{student}\fi
\documentclass[\gkver]{gaokaozhenti}
\begin{document}

\chapter{2011年四川理综}
%% paperType: 理综物理部分

\begin{enumerate}
\item
%% number: 14
%% paperName: 2011年四川理综第 14 题 6 分
%% typeId: 1
%% type: 单选题
%% chapter: 气体性质
%% point: 分子力
%% method: 无
%% score: 6
%% degree: 650
%% duplicateId: 0
%% body:
气体能够充满密闭容器，说明气体分子除相互碰撞的短暂时间外\xzanswer{C}
\fourchoices[answer=C]
{气体分子可以做布朗运动}
{气体分子的动能都一样大}
{相互作用力十分微弱，气体分子可以自由运动}
{相互作用力十分微弱，气体分子间的距离都一样大}

%% answer: C
%% memo:
\memoanswer{
布朗运动是微小颗粒的无规则运动，而不是分子的运动，A 错误；同一温度下，每个气体分子速率不一定相同，其气体分子动能不一定一样大，B 错误；气体分子之间的距离远大于分子力的作用范围，分子力可以认为十分微弱，分子可以自由移动，但分子与分子之间的距离还是不一样大，C 正确，D 错误。
}

\item
%% number: 15
%% paperName: 2011年四川理综第 15 题 6 分
%% typeId: 1
%% type: 单选题
%% chapter: 光学
%% point: 光的电磁说
%% method: 无
%% score: 6
%% degree: 650
%% duplicateId: 0
%% body:
下列说法正确的是\xzanswer{A}
\fourchoices[answer=A]
{甲乙在同一明亮空间，甲从平面镜中看见乙的眼睛时，乙一定能从镜中看见甲的眼睛}
{我们能从某位置通过固定的任意透明介质看见另一侧的所有景物}
{可见光的传播速度总是大于电磁波的传播速度}
{在介质中光总是沿直线传播}

%% answer: A
%% memo:
\memoanswer{
由光路可逆知道 A 正确；由于存在光的全反射，所以不一定从介质看到另一侧的所有景物，B 错误；可见光也是电磁波，它们的速度可以相等，C 错误；只有在同一均匀介质中光才能沿直线传播，D 错误。
}

\item
%% number: 16
%% paperName: 2011年四川理综第 16 题 6 分
%% typeId: 1
%% type: 单选题
%% chapter: 机械波
%% point: 波的图象
%% method: 无
%% score: 6
%% degree: 650
%% duplicateId: 0
%% body:
如图为一列沿 $x$ 轴方向传播的简谐横波在 $t=0$ 时的波形图，当 Q 点在 $t=0$ 时的振动状态传到 P 点时，则\xzanswer{B}
\onepicture[w=6.5cm]{figs/16.png}
\fourchoices[answer=B]
{$1\Ucm<x<3\Ucm$ 范围内的质点正在向 $y$ 轴的负方向运动}
{Q 处的质点此时的加速度沿 $y$ 轴的正方向}
{Q 处的质点此时正在波峰位置}
{Q 处的质点此时运动到 P 处}

%% answer: B
%% memo:
\memoanswer{
由波动图像可知，质点 Q 在 $t=0$ 时刻平衡位置向上振动，所以当 Q 的状态传到 P 时，刚好经历 $3/4$ 周期，Q 处于波谷位置，故 B 正确，C、D 错误；此刻 $1\Ucm<x<3\Ucm$ 范围内的质点均在 $x$ 轴上面，其中 $1\Ucm<x<2\Ucm$ 范围内质点向上运动，$2\Ucm<x<3\Ucm$ 范围内质点向下运动，A 错误。
}

\item
%% number: 17
%% paperName: 2011年四川理综第 17 题 6 分
%% typeId: 1
%% type: 单选题
%% chapter: 曲线运动
%% point: 天体运动
%% method: 无
%% score: 6
%% degree: 450
%% duplicateId: 0
%% body:
据报道，天文学家近日发现了一颗距地球 $40\Uly$ 的“超级地球”，名为“55 Cancrie”，该行星绕母星（中心天体）运行的周期约为地球绕太阳运行周期的 $\frac{1}{480}$，母星的体积约为太阳的 60 倍。假设母星与太阳密度相同，“55 Cancrie”与地球均做匀速圆周运动，则“55 Cancrie”与地球的\xzanswer{B}
\fourchoices[answer=B]
{轨道半径之比约为 $\sqrt[3]{\frac{60}{480}}$}
{轨道半径之比约为 $\sqrt[3]{\frac{60}{480^{2}}}$}
{向心加速度之比约为 $\sqrt[3]{60\times480^{2}}$}
{向心加速度之比约为 $\sqrt[3]{60\times480}$}

%% answer: B
%% memo:
\memoanswer{
母星与太阳密度相等，而体积约为 60 倍，说明母星的质量是太阳的 60 倍。由万有引力提供向心力可得 $G\frac{Mm}{r^{2}}=m\left(\frac{2\pi}{T}\right)^{2}r$，所以 $\frac{M_{\text{母}}}{r_{1}^{3}}\times\frac{r_{2}^{3}}{M_{\text{太}}}=\left(\frac{T_{\text{太}}}{T_{\text{母}}}\right)^{2}$，代入数据得轨道半径之比约为 $\sqrt[3]{60/480^{2}}$，B 正确；由加速度 $a=\left(\frac{2\pi}{T}\right)^{2}r$ 可得，加速度之比为 $\sqrt[3]{60\times480^{4}}$，C、D 错误。
}

\item
%% number: 18
%% paperName: 2011年四川理综第 18 题 6 分
%% typeId: 1
%% type: 单选题
%% chapter: 原子和原子核
%% point: 玻尔模型
%% method: 无
%% score: 6
%% degree: 650
%% duplicateId: 0
%% body:
氢原子从能级 $m$ 跃迁到能级 $n$ 时辐射红光的频率为 $\nu_{1}$，从能级 $n$ 跃迁到能级 $k$ 时吸收紫光的频率为 $\nu_{2}$，已知普朗克常量为 $h$，若氢原子从能级 $k$ 跃迁到能级 $m$，则\xzanswer{D}
\fourchoices[answer=D]
{吸收光子的能量为 $h\nu_{1}+h\nu_{2}$}
{辐射光子的能量为 $h\nu_{1}+h\nu_{2}$}
{吸收光子的能量为 $h\nu_{1}-h\nu_{2}$}
{辐射光子的能量为 $h\nu_{1}-h\nu_{2}$}

%% answer: D
%% memo:
\memoanswer{
由题意可能级 $k$ 处于较高激发态，能级 $m$ 处于较低激发态，而能级 $n$ 处于最低能量状态。所以从能级 $k$ 跃迁到能级 $m$ 时候辐射出光子，能量为 $h\nu_{2}-h\nu_{1}$。
}

\item
%% number: 19
%% paperName: 2011年四川理综第 19 题 6 分
%% typeId: 1
%% type: 单选题
%% chapter: 运动和力
%% point: 超重失重
%% method: 无
%% score: 6
%% degree: 450
%% duplicateId: 0
%% body:
如图是“神舟”系列航天飞船返回舱返回地面的示意图，假定其过程可简化为：打开降落伞一段时间后，整个装置匀速下降，为确保安全着陆，需点燃返回舱的缓冲火箭，在火箭喷气过程中返回舱做减速直线运动，则\xzanswer{A}
\onepicture[h=4.2cm]{figs/19.png}
\fourchoices[answer=A]
{火箭开始喷气瞬间伞绳对返回舱的拉力变小}
{返回舱在喷气过程中减速的主要原因是空气阻力}
{返回舱在喷气过程中所受合外力可能做正功}
{返回舱在喷气过程中处于失重状态}

%% answer: A
%% memo:
\memoanswer{
减速的主要原因是喷出气体，得到反冲作用力，而非空气阻力，B 错误；由动能定理合力做负功，动能减少，减速下降，C 错误；由于减速下降，所以整个舱属于超重状态，D 错误；原来处于平衡状态的返回舱，受到气体向上的反作用力后，绳对返回舱的拉力变小，A 正确。
}

\item
%% number: 20
%% paperName: 2011年四川理综第 20 题 6 分
%% typeId: 2
%% type: 多选题
%% chapter: 交流电
%% point: 交流电
%% method: 无
%% score: 6
%% degree: 450
%% duplicateId: 0
%% body:
如图所示，在匀强磁场中匀速转动的矩形线圈的周期为 $T$，转轴 $O_{1}O_{2}$ 垂直于磁场方向，线圈电阻为 $2\UO$。从线圈平面与磁场方向平行时开始计时，线圈转过 $60^{\circ}$ 时的感应电流为 $1\UA$。那么\xzanswer{AC}
\onepicture[w=4.6cm]{figs/20.png}
\fourchoices[answer=AC]
{线圈消耗的电功率为 $4\UW$}
{线圈中感应电流的有效值为 $2\UA$}
{任意时刻线圈中的感应电动势为 $e=4\cos\frac{2\pi}{T}t\UV$}
{任意时刻穿过线圈的磁通量为 $\Phi=\frac{T}{\pi}\sin\frac{2\pi}{T}t$}

%% answer: AC
%% memo:
\memoanswer{
从图示位置开始计时时，电动势瞬时值满足 $e=NBS\omega\cos\theta=NBS\omega\cos\omega t$，由题意知道，转过 $60^{\circ}$ 时的电动势为 $2\UV$，所以电动势最大值为 $4\UV$，C 选项正确；电流最大值为 $2\UA$，所以有效值为 $\sqrt{2}\UA$，B 错误；$P=I^{2}R=4\UW$，A 选项正确；$e=NBS\omega\cos\theta$，知道 $2=BS\frac{2\pi}{T}\cos60^{\circ}=\frac{\pi}{T}\Phi_{max}$，所以任意时刻的磁通量 $\Phi=\frac{T}{\pi}\sin\frac{2\pi}{T}t$，A、C 正确。
}

\item
%% number: 21
%% paperName: 2011年四川理综第 21 题 6 分
%% typeId: 2
%% type: 多选题
%% chapter: 电场
%% point: 带电粒子的往复运动
%% method: 无
%% score: 6
%% degree: 250
%% duplicateId: 0
%% body:
质量为 $m$ 的带正电小球由空中 A 点无初速度自由下落，在 $t$ 秒末加上竖直向上、范围足够大的匀强电场，再经过 $t$ 秒小球又回到 A 点，不计空气阻力且小球从末落地，则\xzanswer{BD}
\fourchoices[answer=BD]
{整个过程中小球电势能变换了 $\frac{3}{2}mg^{2}t^{2}$}
{整个过程中小球动量增量的大小为 $2mgt$}
{从加电场开始到小球运动到最低点时小球动能变化了 $mg^{2}t^{2}$}
{从 A 点到最低点小球重力势能变化了 $\frac{2}{3}mg^{2}t^{2}$}

%% answer: BD
%% memo:
\memoanswer{
由平均速度关系可知，设下落 $t$ 秒时的速度为 $v$，再次回到 A 点时的速度大小为 $v_{x}$，则满足 $\frac{v}{2}t=\frac{v_{x}+(-v)}{2}t$，即第二次回到 A 点时的速度大小为下落 $t$ 秒时的 2 倍，上升加速度为自由落体加速度的 3 倍，电场力为重力的 4 倍，由动量定理：$4mgt-2mgt=\Delta P$，B 正确；电场力做功对应电势能变化 $W=Fh=4mg\times\frac{1}{2}gt^{2}=2mg^{2}t^{2}$，A 错误；最低点时小球速度为零，所以加电场开始到最低点时，动能变化了 $\Delta E_{k}=mg\times\frac{1}{2}gt^{2}=\frac{1}{2}mg^{2}t^{2}$，C 错误；减速时加速度为自由下落时的 3 倍，所以时间为自由下落的三分之一，总位移为 $h=\frac{1}{2}gt^{2}+\frac{1}{2}\times3g\left(\frac{1}{3}t\right)^{2}=\frac{2}{3}gt^{2}$，所以重力势能的变化 $\Delta E_{p}=\frac{2}{3}mg^{2}t^{2}$，D 正确。
}

\item
%% number: 22
%% paperName: 2011年四川理综第 22 题 10 分
%% typeId: 4
%% type: 实验题
%% chapter: 电路
%% point: 测电动势与内阻
%% method: 无
%% score: 10
%% degree: 450
%% duplicateId: 0
%% body:
\begin{enumerate}
	\item
	（7 分）某研究性学习小组进行了如下实验：如图所示，在一端封闭的光滑细玻璃管中注满清水，水中放一个红蜡做成的小圆柱体 $R$。将玻璃管的开口端用胶塞塞紧后竖直倒置且与 $y$ 轴重合，在 $R$ 从坐标原点以速度 $v_{0}=3\Ucms$ 匀速上浮的同时，玻璃管沿 $x$ 轴正方向做初速为零的匀加速直线运动。同学们测出某时刻 $R$ 的坐标为（4，6），此时 $R$ 的速度大小为\tkanswer[1.2cm]{5} $\Ucms$，$R$ 在上升过程中运动轨迹的示意图是\xzanswer{D}。（$R$ 视为质点）

	\onepicture[align=right, w=5.4cm]{figs/22a.png}

	\fourchoices[answer=D, ispicture=true, h=2.2cm]
	{figs/22b1.png}
	{figs/22b2.png}
	{figs/22b3.png}
	{figs/22b4.png}
	\item
	（10 分）为测量一电源的电动势及内阻
	\begin{steps}
		\item
		在下列三个电压表中选一个改装成量程为 $9\UV$ 的电压表：

		（A）量程为 $1\UV$、内阻大约为 $1\UkO$ 的电压表 $V_{1}$

		（B）量程为 $2\UV$、内阻大约为 $2\UkO$ 的电压表 $V_{2}$

		（C）量程为 $3\UV$、内阻大约为 $3\UkO$ 的电压表 $V_{3}$

		选择电压表\tkanswer[1.0cm]{$V_{3}$}串联\tkanswer[1.0cm]{6}$\UkO$ 的电阻可以改转成量程为 $9\UV$ 的电压表。
		\item
		利用一个电阻箱、一只开关、若开关导线和改装好的电压表（此表用符号 $V_{1}$、$V_{2}$ 或 $V_{3}$ 与一个电阻串联来表示，且可视为理想电压表），在虚线框内画出测量电源电动势及内阻的试验原理示图。

		\onepicture[w=5.2cm]{figs/22c.png}
		\item
		根据以上试验原理电路图进行实验读出电压表示数为 $1.50\UV$ 时、电阻箱值为 $15.0\UO$；电压表示数为 $2.00\UV$ 时、电阻箱的阻值为 $40.0\UO$，则电源的电动势 $E=$\tkanswer[1.0cm]{7.5}$\UV$、内阻 $r=$\tkanswer[1.0cm]{10}$\UO$。
	\end{steps}
\end{enumerate}

%% answer: 
\jdanswer{
\begin{enumerate}
	\item
	$5\Ucms$；D

	\item
	①$V_{3}$，6；②如图所示；③$E=7.5\UV$，$r=10\UO$
\end{enumerate}
}
%% memo:
\memoanswer{
（1）运动时间 $t=2\Us$，所以水平方向平均速度为 $\overline{v}=2\Ucms$，瞬时速度为 $v=4\Ucms$，由速度合成知此刻 $R$ 的速度大小为 $5\Ucms$，由曲线运动条件知道 D 正确。

（2）①由后面第三问可知电压表读数会超过 $2\UV$，所以选电压表 $V_{3}$ 串联一个 $6\UkO$ 的电阻即可改成量程为 $9\UV$ 的电压表。

②如图所示

\onepicture[w=6cm]{figs/22_an.png}

③由电路图知道电压表读数 $1.5\UV$ 时候，路端电压为 $3.5\UV$，读数为 $2\UV$ 时，路端电压为 $6\UV$，由 $E=U+Ir$ 代入数据得到 $E=7.5\UV$，$r=10\UO$。
}

\item
%% number: 23
%% paperName: 2011年四川理综第 23 题 16 分
%% typeId: 6
%% type: 计算题
%% chapter: 运动和力
%% point: 动量定理
%% method: 无
%% score: 16
%% degree: 650
%% duplicateId: 0
%% body:
随着机动车数量的增加，交通安全问题日益凸显。分析交通违法事例，将警示我们遵守交通法规，珍惜生命。一货车严重超载后的总质量为 $49\Ut$，以 $54\Ukmh$ 的速率匀速行驶。发现红灯时司机刹车，货车即做匀减速直线运动，加速度的大小为 $2.5\Umsq$（不超载时则为 $5\Umsq$）。
\begin{enumerate}
	\item
	若前方无阻挡，问从刹车到停下来此货车在超载及不超载时分别前进多远？
	\item
	若超载货车刹车时正前方 $25\Um$ 处停着质量为 $1\Ut$ 的轿车，两车将发生碰撞，设相互作用 $0.1\Us$ 后获得相同速度，问货车对轿车的平均冲力多大？
\end{enumerate}

%% answer: 
\jdanswer{
\begin{enumerate}
	\item
	$45\Um$；$22.5\Um$

	\item
	$9.8\times10^{4}\UN$
\end{enumerate}
}
%% memo:
\memoanswer{
（1）设货车刹车时速度大小为 $v_{0}$、加速度大小为 $a$、末速度为 $v_{t}$、刹车距离为 $s$，$s=\frac{v_{0}^{2}-v_{t}^{2}}{2a}$，代入数据，得超载时 $s_{1}=45\Um$，若不超载 $s_{2}=22.5\Um$。

（2）设货车刹车后经 $s'=25\Um$ 与轿车碰撞时的初速度大小为 $v_{1}$，$v_{1}=\sqrt{v_{0}^{2}-2as'}$，设碰撞后两车共同速度为 $v_{2}$、货车质量为 $M$、轿车质量为 $m$，由动量守恒定律 $Mv_{1}=(M+m)v_{2}$，设货车对轿车的作用时间为 $\Delta t$、平均冲力大小为 $\overline{F}$，由动量定理 $\overline{F}\Delta t=mv_{2}$，代入数据得 $\overline{F}=9.8\times10^{4}\UN$。
}

\item
%% number: 24
%% paperName: 2011年四川理综第 24 题 19 分
%% typeId: 6
%% type: 计算题
%% chapter: 电磁感应
%% point: 双杆问题
%% method: 无
%% score: 19
%% degree: 450
%% duplicateId: 0
%% body:
如图所示，间距 $l=0.3\Um$ 的平行金属导轨 $a_{1}b_{1}c_{1}$ 和 $a_{2}b_{2}c_{2}$ 分别固定在两个竖直面内，在水平面 $a_{1}b_{1}b_{2}a_{2}$ 区域内和倾角 $\theta=37^{\circ}$ 的斜面 $c_{1}b_{1}b_{2}c_{2}$ 区域内分别有磁感应强度 $B_{1}=0.4\UT$、方向竖直向上和 $B_{2}=1\UT$、方向垂直于斜面向上的匀强磁场。电阻 $R=0.3\UO$、质量 $m_{1}=0.1\Ukg$、长为 $l$ 的相同导体杆 K、S、Q 分别放置在导轨上，S 杆的两端固定在 $b_{1}$、$b_{2}$ 点，K、Q 杆可沿导轨无摩擦滑动且始终接触良好。一端系于 K 杆中点的轻绳平行于导轨绕过轻质滑轮自然下垂，绳上穿有质量 $m_{2}=0.05\Ukg$ 的小环。已知小环以 $a=6\Umsq$ 的加速度沿绳下滑，K 杆保持静止，Q 杆在垂直于杆且沿斜面向下的拉力 $F$ 作用下匀速运动。不计导轨电阻和滑轮摩擦，绳不可伸长。取 $g=10\Umsq$，$\sin37^{\circ}=0.6$，$\cos37^{\circ}=0.8$。求：
\begin{enumerate}
	\item
	小环所受摩擦力的大小；
	\item
	Q 杆所受拉力的瞬时功率。
\end{enumerate}
\onepicture[align=right, w=6.8cm]{figs/24.png}

%% answer: 
\jdanswer{
\begin{enumerate}
	\item
	$0.2\UN$

	\item
	$2\UW$
\end{enumerate}
}
%% memo:
\memoanswer{
（1）设小环受到的摩擦力大小为 $F_{f}$，由牛顿第二定律，有 $m_{2}g-F_{f}=m_{2}a$，代入数据，得 $F_{f}=0.2\UN$。

（2）设通过 K 杆的电流为 $I_{1}$，K 杆受力平衡，有 $F_{f}=B_{1}I_{1}l$，设回路总电流为 $I$，总电阻为 $R_{\text{总}}$，有

$I=2I_{1}$，$R_{\text{总}}=\frac{3}{2}R$，

设 Q 杆下滑速度大小为 $v$，产生的感应电动势为 $E$，有

$I=\frac{E}{R_{\text{总}}}$，$E=B_{2}lv$，$F+m_{1}g\sin\theta=B_{2}Il$，拉力的瞬时功率为 $P=Fv$，联立以上方程得到 $P=2\UW$。
}

\item
%% number: 25
%% paperName: 2011年四川理综第 25 题 20 分
%% typeId: 6
%% type: 计算题
%% chapter: 磁场
%% point: 带电粒子在复合场中的运动
%% method: 无
%% score: 20
%% degree: 250
%% duplicateId: 0
%% body:
如图所示：正方形绝缘光滑水平台面 WXYZ 边长 $l=1.8\Um$，距地面 $h=0.8\Um$。平行板电容器的极板 CD 间距 $d=0.1\Um$ 且垂直放置于台面，C 板位于边界 WX 上，D 板与边界 WZ 相交处有一小孔。电容器外的台面区域内有磁感应强度 $B=1\UT$、方向竖直向上的匀强磁场。电荷量 $q=5\times10^{-13}\UC$ 的微粒静止于 W 处，在 CD 间加上恒定电压 $U=2.5\UV$，板间微粒经电场加速后由 D 板所开小孔进入磁场（微粒始终不与极板接触），然后由 XY 边界离开台面。在微粒离开台面瞬时，静止于 X 正下方水平地面上 A 点的滑块获得一水平速度，在微粒落地时恰好与之相遇。假定微粒在真空中运动、极板间电场视为匀强电场，滑块视为质点，滑块与地面间的动摩擦因数 $\mu=0.2$，取 $g=10\Umsq$。
\begin{enumerate}
	\item
	求微粒在极板间所受电场力的大小并说明两板的极性；
	\item
	求由 XY 边界离开台面的微粒的质量范围；
	\item
	若微粒质量 $m_{0}=1\times10^{-13}\Ukg$，求滑块开始运动时所获得的速度。
\end{enumerate}
\onepicture[align=right, w=7.2cm]{figs/25.png}

%% answer: 
\jdanswer{
\begin{enumerate}
	\item
	$1.25\times10^{-11}\UN$；C 板为正极，D 板为负极。

	\item
	$8.1\times10^{-14}\Ukg<m\leqslant2.89\times10^{-13}\Ukg$

	\item
	$4.15\Ums$
\end{enumerate}
}
%% memo:
\memoanswer{
（1）微粒在极板间所受电场力大小为 $F=\frac{qU}{d}$，代入数据 $F=1.25\times10^{-11}\UN$，由微粒在磁场中的运动可判断微粒带正电荷，微粒由极板间电场加速，故 C 板为正极，D 板为负极。

（2）若粒子的质量为 $m$，刚进入磁场时的速度大小为 $v$，由动能定理 $qU=\frac{1}{2}mv^{2}$，微粒在磁场中做匀速圆周运动，洛伦磁力充当向心力，若圆周运动半径为 $R$，有 $qvB=m\frac{v^{2}}{R}$，微粒要从 XY 边界离开台面，则圆周运动的边缘轨迹如图所示，半径的极小值与极大值为 $R_{1}=\frac{l}{2}$，$R_{2}=l-d$，代入数据，有 $8.1\times10^{-14}\Ukg<m\leqslant2.89\times10^{-13}\Ukg$。

\onepicture[w=5.4cm]{figs/25_an1.png}
\onepicture[w=5.2cm]{figs/25_an2.png}

（3）如图，微粒在台面以速度为 $v$ 做以 O 点位圆心，$R$ 为半径的圆周运动；从台面边缘 P 点沿与 XY 边界成 $\theta$ 角飞出做平抛运动，落地点 Q 点，水平位移 $s$，下落时间 $t$。设滑块质量为 $M$，滑块获得的速度 $v_{0}$ 后在 $t$ 内与平台前侧面成 $\varphi$ 角度方向，以加速度 $a$ 做匀减速直线运动到 Q，经过位移为 $k$。由几何关系，可得 $\cos\theta=\frac{l-R}{R}$，根据平抛运动，$t=\sqrt{\frac{2h}{g}}$，$s=vt$，对于滑块，由牛顿定律及运动学方程，有 $\mu Mg=Ma$，$k=v_{0}t-\frac{1}{2}at^{2}$，再有余弦定理，$k^{2}=s^{2}+(d+R\sin\theta)^{2}-2s(d+R\sin\theta)\cos\theta$，及正弦定理，$\frac{\sin\varphi}{s}=\frac{\sin\theta}{k}$，联立，代入数据，解得：$v_{0}=4.15\Ums$，$\varphi=\arcsin0.8$（或 $\varphi=53^{\circ}$）。

\onepicture[w=11cm]{figs/25_an3.png}
}

\end{enumerate}

\end{document}
